\section{Global complexity: a precise remaining target}
\label{sec:global}

\subsection{The family-query bound removes one parameter explosion}

Suppose a represented stage previously tried up to $m^k$ phase vectors,
and the only information requested from that family is feasibility and
the minimum number of cyclic covering components. For the affine model
of this article, compilation and optimization have polynomial bit cost
in the listed matrix dimensions and input integer lengths. Thus the
$m^k$ term can genuinely be removed from that stage. The emitted witness
and its independent verification are polynomial as well. This conclusion
does not require $k=O(\log n)$ or $m$ to be a small integer.

There is no corresponding conclusion for a stage which must distinguish
all seam assignments under future attachments, require many selected
boundary preimages to be connected, or test unknown normal-surface
embeddings. Sections~\ref{sec:surface-families} and \ref{sec:boundaries}
give examples and formal barriers to those extensions.

\subsection{A composition theorem with every hypothesis exposed}

\begin{theorem}[Conditional represented-algorithm bound]
\label{thm:conditional-global}
Consider a complete unknot-recognition algorithm on $n$-crossing diagrams.
Suppose that, on every execution:
\begin{enumerate}
\item the number of represented states and transitions, including all
failed branches and resets, is at most $\exp(C\log^2(n+2))$;
\item every listed state and certificate has bit size at most
$\exp(C\log^2(n+2))$;
\item every transition is either a verified operation proved polynomial
in that bit size, or another routine whose total cost has the same
quasi-polynomial bound;
\item the representation and transition relation are sound and complete
for the required topological decision.
\end{enumerate}
Then the total bit time, including certificate checks and output, is
$\exp(O(\log^2(n+2)))$.
\end{theorem}
\begin{proof}
A fixed polynomial in $\exp(C\log^2(n+2))$ is
$\exp(O(\log^2(n+2)))$. Multiplying this by the assumed bound on the
number of transitions preserves the form. Reading and emitting every
state is included in the same product. The soundness and completeness
hypothesis is needed to identify the resulting algorithm as a complete
recognizer rather than an exact subroutine or a partial search.
\end{proof}

This bookkeeping theorem is elementary. Its usefulness is diagnostic:
the arithmetic optimizer and normal-coordinate extractor establish
particular transition-cost statements in item (3), under their explicit
input contracts. Scalar-space inheritance avoids repeated work on an
already represented complex. None establishes items (1), (2), or (4) for
unrestricted knots. In particular, the theorem is not presented as a
new solution to the central complexity problem.

\subsection{Why a termination potential is insufficient}

A lexicographically decreasing list with $L$ entries each in
$\{0,\ldots,g\}$ can visit $(g+1)^L$ values. The corresponding integer
potential is
\[
 P(a_1,\ldots,a_L)=\sum_{i=1}^L a_i(g+1)^{L-i}.
\]
Decreasing an earlier coordinate and resetting all later coordinates to
$g$ decreases $P$, but can make progress of only one unit. This is exactly
the type of accounting behind the hierarchy iteration bound
$L(g+1)^L$ in Lackenby's Proposition 9.1~\cite{lackenby2026}.

A sufficient numerical regime for this particular bound is
\[
 \log(L+1)+L\log(g+1)=O(\log^2(n+2)),
\]
together with control of the per-step bit cost. The first term also
covers $g=0$, when the bound reduces to $L$. A merely linear estimate
for $L$ does not supply it. Faster gcds, shorter certificates, or reduced
commutant solves do not change the number of possible resets. A new
global argument would need, for example, a stronger amortized progress
measure, a decomposition theorem limiting genuinely new states, or
an operation that processes a whole relevant reset family symbolically.

\subsection{The two geometric interfaces must not be conflated}

The normal extractor produces partial signed interval maps, with source
and target stacks and explicit endpoint data. The family optimizer
receives total translations of a common cyclic fibre and affine
parameter constraints. These are different mathematical objects. A
sound bridge could identify a cyclic block inside the interval data,
certify that all its possible completions are represented, and pass the
induced congruences to the optimizer. No such general bridge is included.

Likewise, computing the local complement cells is less than constructing
their full residual handle structure, including every attachment and
boundary pattern. General orbit counting can determine connected pieces,
but subsequent genus, orientation, and vertical-annulus data must retain
the appropriate provenance. Theorem 14.4 of the geometric literature
identifies the desired full uniform-type handle-structure interface~\cite{lackenby2026}; the present
implementation supplies a concrete first portion of it.

\subsection{The algebraic route retains its intrinsic size barrier}

Transporting an endomorphism space can avoid repeatedly solving for the
same information. A certified local scalar endomorphism algebra can stop
a futile search for scalar direct summands. Neither fact bounds the
number of surviving homological generators or the size of the parent
commutant. A complex with a local endomorphism algebra can still be large.
The first computation of that algebra may dominate, and a bounded search
can remain inconclusive when no dimension-saturated generator is found.

The measured implementation therefore retains caps, transactional failure
behavior, and separate optional controls. A local algebra certificate
means ``no nontrivial splitting by the admitted scalar maps.'' It does
not mean ``unknot,'' ``small homology,'' or ``no more general
continuation-compatible compression.''
